\section{Related Work}
\label{sec:related}

We group prior work by the pillar of the problem it addresses and close with a comparison
(Table~\ref{tab:compare}).

\subsection{Memory forensics and AI-assisted triage}
Volatility is the standard open framework for reconstructing operating-system state from
RAM; its \texttt{malfind} plugin flags executable, writable memory not backed by a file, a classic sign of
injected code, but also fires on just-in-time compilers and runtime stubs. Carrier et al.~\cite{carrier2022detecting}
engineer features from memory dumps and detect obfuscated malware with 99\% accuracy, but the output is a
label with no account of the cryptography involved. volGPT~\cite{oh2024volgpt} passes Volatility artefacts to
an LLM to rank ransomware processes, and ForensicLLM~\cite{sharma2025forensicllm} shows that a fine-tuned local
model can answer forensic questions as well as larger cloud models. Scanlon et al.~\cite{scanlon2023chatgpt}
find GPT-4 useful for some forensic tasks but stress that the analyst must catch its mistakes. None of these
asks whether planted evidence can bias the model, and none requires the model to cite evidence.

\subsection{Keys and ciphers in memory and binaries}
Halderman et al.\ (2008) showed that AES keys can be found in a memory image by searching for
byte strings whose expansion matches the following key schedule, the basis of our AES key finder.
Fernandez de Loaysa Babiano et al.~\cite{babiano2023evaluation} recover over 90\% of Salsa20 key and nonce
pairs from the memory of running ransomware, and Baier et al.~\cite{baier2024tls} systematise how TLS key
material is located in memory. These works each target one key type in one setting. Casino et
al.~\cite{casino2025not} classify high-entropy file segments to separate encryption from compression, without
naming the algorithm. Meijer et al.~\cite{meijer2021wheres} identify known and proprietary primitives in binary
code with data-flow graphs, and Barnes and Ghafarian~\cite{barnes2026machine} detect ransomware from PE-header
features; both operate on a single binary rather than a whole memory image. Entropy as a packing or
encryption indicator goes back to Lyda and Hamrock (2007), whose measured averages (7.18 bits per
byte for encrypted executables) set our threshold.

\subsection{Cipher cost and weak keys}
Baig et al.~\cite{baig2024comparative} measure the speed and memory
cost of AES, 3DES and RSA, which explains why malware uses symmetric ciphers for bulk data and public
keys only to wrap them~\cite{oz2022survey}. For RSA, Heninger et al.\ (2012) found widespread
weak keys in network devices by computing pairwise GCDs, Pelofske~\cite{pelofske2024efficient} speeds up that
batch GCD, and Ruzai et al.~\cite{ruzai2024concurrent} factor moduli linked by weak key equations. These operate
on key collections or in theory; strength is judged against NIST SP~800-57 and the transition
rules of SP~800-131A, with the 829-bit RSA-250 factorisation as the public
record that marks a modulus as broken.

\subsection{Tokens and credentials}
Raghu~\cite{raghu2025context} quantifies the cost of hardening JWT and OAuth in production, and Rujichaikul and
Rassameeroj~\cite{rujichaikul2025token} detect token misuse with rule-based signatures at 81.4\% accuracy.
RFC~8725 and RFC~4226 define what makes a JWT or an OTP secret weak. Neither line of
work looks for credentials left in memory.

\subsection{LLM robustness to injected evidence}
Greshake et al.~\cite{greshake2023not} introduced indirect prompt injection through data an application
retrieves, which is exactly the forensic setting. Spotlighting~\cite{hines2024spotlighting} marks untrusted
input with delimiters or encodings and cuts attack success substantially;
Liu et al.~\cite{liu2024formalizing} benchmark attacks and defences and
conclude that no single defence is sufficient. None of these is evaluated on forensic evidence or combined
with a deterministic check of the model's claims and numbers.

\subsection{Gap}
Table~\ref{tab:compare} summarises the comparison. TriageGuard is, to our knowledge, the first system to
combine memory artefacts, cipher identification with an evidence basis, recovery and NIST assessment of keys
carved from the image, and an LLM report whose exposure to injection is both measured and defended.

\textit{Why Table~\ref{tab:compare} has these columns and rows.} The six columns are the six capabilities our
problem statement needs, in pipeline order: reading a memory image, identifying the cipher, recovering keys,
judging key strength, writing a report with an LLM, and testing that report against injected evidence. A system
missing any one of them cannot answer the question in Section~\ref{sec:intro}. The rows are the closest work for
each capability from our 20-paper survey: for each column we kept the paper whose method is nearest to ours, so
the table shows where each pillar is already solved and that no single work covers all six.

\begin{table}[t]
\centering
\caption{Comparison with representative prior work}
\label{tab:compare}
\footnotesize
\setlength{\tabcolsep}{3pt}
\begin{tabular}{@{}lcccccc@{}}
\toprule
Work & \rotatebox{90}{Memory image} & \rotatebox{90}{Cipher ID} & \rotatebox{90}{Key recovery} & \rotatebox{90}{Weak-key test} & \rotatebox{90}{LLM report} & \rotatebox{90}{Injection tested}\\
\midrule
Carrier et al.~\cite{carrier2022detecting} & \done & & & & & \\
volGPT~\cite{oh2024volgpt} & \done & & & & \done & \\
ForensicLLM~\cite{sharma2025forensicllm} & & & & & \done & \\
Babiano et al.~\cite{babiano2023evaluation} & \done & & \done & & & \\
Meijer et al.~\cite{meijer2021wheres} & & \done & & & & \\
Spotlighting~\cite{hines2024spotlighting} & & & & & \done & \done\\
\midrule
\textbf{TriageGuard} & \done & \done & \done & \done & \done & \done\\
\bottomrule
\end{tabular}
\end{table}
